% Introduction of the Workbook (model: WordPerfect 5.1 Workbook pp. v-viii:
% before you start, the lesson files, your computer, the screens, the default
% settings, how the lessons are grouped; docs/pubs/originals/
% wordperfect-51.md s5).  The model's LEARN files and printer definition
% have no InitechWord counterpart: the lessons make their own documents.
% GATES: test-pubs-iword
\Gates{test-pubs-iword}
\FrontChapter{Introduction}
\markboth{Introduction}{Introduction}
This workbook teaches you InitechWord by having you use it.  Each lesson
gives you an office document to produce -- a note, a memo, a list of
procedures, a letter -- and takes you through it keystroke by keystroke.
Do the lessons at your computer, with InitechWord running.

\WPSide{Before You Start}
Read \textit{Getting Started} in the \textit{InitechWord Reference} first.
It tells you how InitechWord is supplied, how to start it, and how its keys
and menus work.  This workbook assumes you can start InitechWord from the
FLAIR desktop and know where the function keys are.  Lay the function-key
template above the function keys of your keyboard.

\WPSide{The Documents You Make}
The lessons make their own documents: you type each one, and save it on the
data volume under the name the lesson gives.  Several lessons go on with a
document an earlier lesson saved: Lesson~3 retrieves the memo of Lesson~2,
for example.  Do those lessons in order, and use the names the lessons give,
so that each document is where the next lesson expects it.

\WPSide{The Screens in This Workbook}
The panels headed \textit{Compare your screen} show what InitechWord's
editing screen shows at that point of the lesson, cell for cell: 76
characters across and 24 lines down, with the status line at the bottom.
The lines below the last one that holds something are left out; the status
line is always printed.  The cursor is shown as a black cell, and a block as a grey band.
Bold, italic and underlined text is printed as InitechWord draws it.

If your screen differs, look back over the steps since the last panel: a
key pressed twice or a word typed differently is the usual cause.  You can
always clear the screen (press Exit, F7, then type \Typed{n} and \Typed{n})
and start the lesson again.

\WPSide{The Settings}
The lessons assume the settings InitechWord gives a new document: one inch
margins, single spacing, tab stops every half inch, and 66 characters on a
line.  A lesson that changes a setting says so.

\WPSide{How the Lessons Are Grouped}
The lessons are in three sections:
\begin{description}
\item[Fundamentals I] Lessons 1 to 8: typing, saving, retrieving, moving the
cursor, bold and underline, codes, indents and pages.  Do these in order.
\item[Fundamentals II] Lessons 9 to 16: blocks, moving and copying text,
deleting, and the appearance and size of text.  Lessons 9 to 13 go on with
documents of the first section.
\item[Special Features] Lessons 17 to 30: margins, justification, line
spacing, combining documents, the messages and a set of techniques to dip
into (17 to 22), then complete documents that use the features together: a
letter, minutes, a final copy, a form, a revised list, a notice, a table,
and text copied between documents (23 to 30).
\end{description}
Each lesson ends by clearing the screen or leaving InitechWord, so you can
stop after any lesson.  Each begins with InitechWord just started, its
screen clear.

\WPSide{Where to Learn More}
Each lesson names the features it teaches.  To learn more about a feature,
read its entry in the \textit{InitechWord Reference}; Appendix A of this
workbook lists every feature with the lessons that use it.
